
\section{与镜像对称的关系}
\label{sec:1-4}

到目前为止，我们已经说明了弗洛尔理论对辛拓扑问题的各种应用。另一条应用方向，是证明孔采维奇在约 1994 年提出的著名同调镜像对称猜想\textbf{[Kon93,95I]}；第一署名作者引入的 $A_\infty$ 范畴\liftIndex{I006}\textbf{[Fuk93]}，在这一猜想的表述中起着关键作用。（瓦法约在 2000 年把这个构想解释成物理学语言\textbf{[Vaf99]}。）要实现这一目标，还有许多事情有待完成。这里我们想简要说明本书工作与同调镜像对称的关系（其他说明另见\textbf{[Fuk02I,02III,03I,03II]}）。

为此，回忆定理 B 后紧接着的一段所引入的集合
\[
  \mathcal M(L)
  =\bigl(\operatorname{Zero}(\pi_{\mathrm{amb}})
    \cap\operatorname{Zero}(\mathfrak{PO})\bigr)/{\sim},
\]
它由定界余链的规范等价类组成。（在本节此处及以下，均假设 $L$ 是 $M$ 的相对 spin 拉格朗日子流形。）由于这个集合在同调镜像对称的应用中居于核心地位，我们先陈述一个把它描述成形式概形\liftIndex{I017}的定理。设 $C_0,C_1$ 是两个有限生成自由 $\Lambda_{0,\mathrm{nov}}$-模。分别选取 $C_0$ 的基 $\{\mathbf v_1,\ldots,\mathbf v_m\}$ 和 $C_1$ 的基 $\{\mathbf w_1,\ldots,\mathbf w_\ell\}$。所谓\emph{形式映射}\liftIndex{I018} $\mathfrak P:C_0\to C_1$，是指线性组合 $\mathfrak P=\sum_j\mathfrak P_j\mathbf w_j$，其中每个 $\mathfrak P_j$ 都是形式幂级数
\begin{equation}
  \mathfrak P_j\!\left(\sum_{i=1}^{m}s_i\mathbf v_i\right)
  =\sum_{k_1=0}^{\infty}\cdots\sum_{k_m=0}^{\infty}
    a^j_{k_1,\ldots,k_m}s_1^{k_1}\cdots s_m^{k_m}.
  \tag{1.15}\label{eq:1-15}
\end{equation}
这里的系数 $a^j_{k_1,\ldots,k_m}\in\Lambda_{0,\mathrm{nov}}$。注意，一般而言，如果把 $\Lambda_{0,\mathrm{nov}}$ 中的元素代入 $s_j$，公式~\eqref{eq:1-15} 未必有意义；但若代入 $\Lambda^+_{0,\mathrm{nov}}$ 中的元素，它便有意义。这是因为在后一种情形，\eqref{eq:1-15} 的右端在 $\Lambda^+_{0,\mathrm{nov}}$-进拓扑下收敛。因此可以令
\[
  \bigl(\mathfrak P^{-1}(0)\bigr)(\Lambda^+_{0,\mathrm{nov}})
  :=\left\{
    \left.\sum_j s_j\mathbf v_j\in\Lambda^+_{0,\mathrm{nov}}C_0
    \ \right|\
    \mathfrak P\!\left(\sum_{i=1}^{m}s_i\mathbf v_i\right)=0
  \right\},
\]
并把这个集合看作仿射形式概形
$\operatorname{Spf}(R(\mathfrak P;\Lambda_{0,\mathrm{nov}}))$
的 $\Lambda^+_{0,\mathrm{nov}}$-值点（或刚性点）；该形式概形对应于（交换、非诺特的）完备赋值环
\[
  R(\mathfrak P;\Lambda_{0,\mathrm{nov}})
  =\frac{\Lambda_{0,\mathrm{nov}}[[s_1,\ldots,s_m]]}
  {(\mathfrak P_1,\ldots,\mathfrak P_\ell)}.
\]
这里 $(\mathfrak P_1,\ldots,\mathfrak P_\ell)$ 是形式幂级数环
$\Lambda_{0,\mathrm{nov}}[[s_1,\ldots,s_m]]$
中由 $\mathfrak P_1,\ldots,\mathfrak P_\ell$ 生成的理想。换言之，有如下等同：
\[
  \bigl(\mathfrak P^{-1}(0)\bigr)(\Lambda^+_{0,\mathrm{nov}})
  =\operatorname{Hom}\bigl(R(\mathfrak P;\Lambda_{0,\mathrm{nov}}),
    \Lambda_{0,\mathrm{nov}}\bigr),
\]
其中右端是所有连续 $\Lambda_{0,\mathrm{nov}}$-代数同态的集合（见第 5.4.1 小节的命题 5.4.15）。

对每个 $\mathfrak b\in H^2(M;\Lambda^+_{0,\mathrm{nov}})$，考虑映射
\[
  \pi_{\mathrm{amb}}:\mathcal M_{\mathrm{weak,def}}(L)
  \longrightarrow H^2(M;\Lambda^+_{0,\mathrm{nov}})
\]
在 $\mathfrak b$ 上的原像 $\pi_{\mathrm{amb}}^{-1}(\mathfrak b)$，并定义
\[
  \mathcal M^{\mathfrak b}_{\mathrm{weak}}(L)
  =\pi_{\mathrm{amb}}^{-1}(\mathfrak b),
  \qquad
  \mathcal M^{\mathfrak b}(L)
  =\pi_{\mathrm{amb}}^{-1}(\mathfrak b)
   \cap\operatorname{Zero}(\mathfrak{PO}).
\]

\begin{theorem}{M}\label{thm:intro-M}
存在映射
\[
  \mathfrak P^{\mathfrak b}_{L,\mathrm{weak}}:
  H^1(L;\Lambda^+_{0,\mathrm{nov}})
  \longrightarrow H^2(L;\Lambda^+_{0,\mathrm{nov}})
\]
以及 $\bigl((\mathfrak P^{\mathfrak b}_{L,\mathrm{weak}})^{-1}(0)\bigr)
(\Lambda^+_{0,\mathrm{nov}})$ 上的一个等价关系 $\sim$，使得
\[
  \mathcal M^{\mathfrak b}_{\mathrm{weak}}(L)
  =\bigl((\mathfrak P^{\mathfrak b}_{L,\mathrm{weak}})^{-1}(0)
    (\Lambda^+_{0,\mathrm{nov}})\bigr)/{\sim}.
\]
作为 $\Lambda_{0,\mathrm{nov}}$-代数，
$R(\mathfrak P^{\mathfrak b}_{L,\mathrm{weak}};\Lambda_{0,\mathrm{nov}})$
的连续同构类，在辛微分同胚意义下只依赖于对子 $(M,L)$。若改用
\[
  \mathfrak P^{\mathfrak b}_L:H^1(L;\Lambda^+_{0,\mathrm{nov}})
  \longrightarrow H^2(L;\Lambda^+_{0,\mathrm{nov}}),
\]
则对 $\mathcal M^{\mathfrak b}(L)$ 也有类似结论。
\end{theorem}

定理 M 是定理 5.4.12 和命题 5.4.15 的推论，其证明见第 5.4.6 小节。

现在考虑一对相对 spin 的拉格朗日子流形 $(L_0,L_1)$。给定
$\mathfrak b\in H^2(M;\Lambda^+_{0,\mathrm{nov}})$ 和
$b_i\in\mathcal M^{\mathfrak b}_{\mathrm{weak}}(L_i)$，借助定理 M 中的等同，以
\[
  \varphi_{b_i}:R(\mathfrak P^{\mathfrak b}_{L_i};\Lambda_{0,\mathrm{nov}})
  \longrightarrow\Lambda_{0,\mathrm{nov}}
\]
表示与 $b_i$ 的一个代表相应的同态。当对子 $(L_1,L_0)$ 相对 spin 时，有如下结论。

\begin{theorem}{N}\label{thm:intro-N}
存在一个复形
$\bigl(\mathfrak D(L_1,L_0;\Lambda_{0,\mathrm{nov}}),\mathfrak d\bigr)$，其中
$\Lambda_{0,\mathrm{nov}}$-模
$\mathfrak D(L_1,L_0;\Lambda_{0,\mathrm{nov}})$
是环对
\[
  \bigl(R(\mathfrak P^{\mathfrak b}_{L_1};\Lambda_{0,\mathrm{nov}}),
  R(\mathfrak P^{\mathfrak b}_{L_0};\Lambda_{0,\mathrm{nov}})\bigr)
\]
上的有限生成完备双模，而上边界映射
\[
  \mathfrak d:\mathfrak D(L_1,L_0;\Lambda_{0,\mathrm{nov}})
  \longrightarrow\mathfrak D(L_1,L_0;\Lambda_{0,\mathrm{nov}})
\]
关于滤过在 $\mathfrak D(L_1,L_0;\Lambda_{0,\mathrm{nov}})$ 上诱导的拓扑连续。

在同构意义下，这个复形与所作的各种选择无关，并具有下列性质。

\par\smallskip\noindent\textup{(N.1)}\quad
取一对定界余链 $b_i\in\mathcal M^{\mathfrak b}(L_i)$（$i=0,1$），并分别以
\[
  \varphi_{b_i}:R(\mathfrak P^{\mathfrak b}_{L_i};\Lambda_{0,\mathrm{nov}})
  \longrightarrow\Lambda_{0,\mathrm{nov}}
\]
表示其代表。则 $\Lambda_{0,\mathrm{nov}}$-双模复形
\[
\begin{split}
\bigl(&{\Lambda_{0,\mathrm{nov}}}_{\varphi_{b_1}}
  \otimes_{R(\mathfrak P^{\mathfrak b}_{L_1};\Lambda_{0,\mathrm{nov}})}
  \mathfrak D(L_1,L_0;\Lambda_{0,\mathrm{nov}})
  \otimes_{R(\mathfrak P^{\mathfrak b}_{L_0};\Lambda_{0,\mathrm{nov}})}
  {_{\varphi_{b_0}}\Lambda_{0,\mathrm{nov}}},\\
&1\otimes\mathfrak d\otimes1\bigr)
\end{split}
\]
的上同调，作为 $\Lambda_{0,\mathrm{nov}}$-双模，同构于弗洛尔上同调
\[
  HF\bigl((L_1,\mathbf b_1),(L_0,\mathbf b_0);\Lambda_{0,\mathrm{nov}}\bigr).
\]
这里 $\mathbf b_i=(\mathfrak b,b_i)$（$i=0,1$）。

\par\smallskip\noindent\textup{(N.2)}\quad
若 $\phi_i:(M,L_i)\to(M,L_i')$（$i=0,1$）是哈密顿同痕，则存在链同伦等价
\[
\begin{split}
  (\phi_1,\phi_0)_*:\;&
  \bigl(\mathfrak D(L_1,L_0;\Lambda_{0,\mathrm{nov}}),\mathfrak d\bigr)
  \otimes_{\Lambda_{0,\mathrm{nov}}}\Lambda_{\mathrm{nov}}\\
  &\longrightarrow
  \bigl(\mathfrak D(L_1',L_0';\Lambda_{0,\mathrm{nov}}),\mathfrak d\bigr)
  \otimes_{\Lambda_{0,\mathrm{nov}}}\Lambda_{\mathrm{nov}},
\end{split}
\]
它位于环同构
\[
  \phi_{i*}:R(\mathfrak P^{\mathfrak b}_{L_i};\Lambda_{0,\mathrm{nov}})
  \longrightarrow
  R(\mathfrak P^{\mathfrak b}_{L_i'};\Lambda_{0,\mathrm{nov}})
\]
之上。在 (N.2) 给出的同构下，$(\phi_1,\phi_0)_*$ 诱导 (G.4) 所给出的弗洛尔上同调同构。\translatornote{原书此处写作 (G.5)，但定理 G 中哈密顿同痕诱导的弗洛尔上同调同构编号为 (G.4)；(G.5) 是双模结构陈述，故据上下文改正。}
\end{theorem}

上面的 (N.1) 还可以作如下重新解释：由形式概形的一般理论，复形\translatornote{原书此处写作 $\mathfrak D(L'_1,L'_0)$，但本段是在重述 (N.1)，其底形式概形及随后纤维均使用 $L_1,L_0$；故改为无撇记号。}
$\bigl(\mathfrak D(L_1,L_0;\Lambda_{0,\mathrm{nov}}),\mathfrak d\bigr)$
在仿射形式概形
\[
  \operatorname{Spf}\bigl(R(\mathfrak P^{\mathfrak b}_{L_1};\Lambda_{0,\mathrm{nov}})\bigr)
  \times
  \operatorname{Spf}\bigl(R(\mathfrak P^{\mathfrak b}_{L_0};\Lambda_{0,\mathrm{nov}})\bigr)
\]
上定义了凝聚层导出范畴中的一个对象。它在（刚性）点 $(b_1,b_0)$ 处的纤维就是弗洛尔上同调。定理 N 的证明见第 5.4.6 小节。

\begin{remark}{1.16}\label{rem:1-16}
定理 N 中考虑的是 $\mathcal M^{\mathfrak b}(L)$。也可以改用
$\mathcal M^{\mathfrak b}_{\mathrm{weak}}(L)$ 而得到类似结论。具体说，利用态射
\[
  \mathfrak{PO}^{\mathfrak b}_L:
  \mathcal M^{\mathfrak b}_{\mathrm{weak}}(L)
  \longrightarrow\Lambda^+_{0,\mathrm{nov}},
\]
它对应于连续环同态
\[
  \Lambda^+_{0,\mathrm{nov}}
  \longrightarrow
  R(\mathfrak P^{\mathfrak b}_{L,\mathrm{weak}};\Lambda_{0,\mathrm{nov}}),
\]
可以定义纤维积
\[
  R(\mathfrak P^{\mathfrak b}_{L_1,\mathrm{weak}};\Lambda_{0,\mathrm{nov}})
  \mathbin{\widehat\otimes}_{\Lambda^+_{0,\mathrm{nov}}}
  R(\mathfrak P^{\mathfrak b}_{L_0,\mathrm{weak}};\Lambda_{0,\mathrm{nov}}).
\]
于是 $\mathfrak D(L_1,L_0;\Lambda_{0,\mathrm{nov}})$ 可以扩张成上述环上的模所成的链复形。它在刚性点处的纤维给出一族弗洛尔上同调，其参数空间为
\[
  \mathcal M^{\mathfrak b}_{\mathrm{weak}}(L_1)
  \mathbin{\times}_{\Lambda^+_{0,\mathrm{nov}}}
  \mathcal M^{\mathfrak b}_{\mathrm{weak}}(L_0),
\]
其中纤维积是相对于 $\mathfrak{PO}^{\mathfrak b}_{L_i}$ 取得的。
\end{remark}

本节其余部分讨论本书所得结果与镜像对称的关系。除去以后偶尔作类似说明的需要，这些讨论不会用于本书其余部分。

在镜像对称的应用中，最相关的一类拉格朗日子流形 $L\subset M$ 是这样的：$M$ 为卡拉比--丘流形，而 $L$ 的马斯洛夫类为零。因此在下面的讨论中，我们将限于这一类拉格朗日子流形。（不过，本书结果也适用于法诺情形；它在镜像对称中同样起重要作用。）

首先注意，对这样的拉格朗日子流形，不需要在 $\Lambda_{0,\mathrm{nov}}$ 中使用记录圆盘马斯洛夫指标的形式参数 $e$（见第 5.4.6 小节）。因此我们分别用较小的环 $\Lambda_0,\Lambda,\Lambda_+$ 取代 $\Lambda_{0,\mathrm{nov}},\Lambda_{\mathrm{nov}},\Lambda^+_{0,\mathrm{nov}}$（见第 1 章末尾的 (Conv.5)）。注意，此时 $\Lambda$ 是一个域。于是可以在 $H(L;\Lambda_0)$ 上定义滤过 $A_\infty$ 代数结构；也可以像对 $\Lambda^+_{0,\mathrm{nov}}$ 那样，以相同方式在这个较小的环上定义形式概形。为了明确表示这一改变，我们把相应的形式概形记为 $\mathcal M(L;\Lambda_+)$，以强调它对 $\Lambda_+$ 的依赖。下面这个定理 G 和定理 M 的版本，适合用来描述本书所研究的弗洛尔上同调与镜像对称之间的关系。

\begin{corollary}{O}\label{cor:intro-O}
设 $(M,\omega)$ 满足 $c_1(M)=0$。设 $L$ 是 $M$ 中马斯洛夫类为零的相对 spin 拉格朗日子流形，并且
$\mathfrak b\in H^2(M;\Lambda_+)$。则存在映射
\[
  \mathfrak P^{\mathfrak b}_L:H^1(L;\Lambda_+)
  \longrightarrow H^2(L;\Lambda_+),
\]
使得集合
\[
  \mathcal M^{\mathfrak b}(L;\Lambda_+)
  =(\mathfrak P^{\mathfrak b}_L)^{-1}(0)(\Lambda_+)
  =\{b\in H^1(L;\Lambda_+)\mid\mathfrak P^{\mathfrak b}_L(b)=0\}
\]
具有下列性质。设 $(L_0,L_1)$ 是 $M$ 中一对这样的拉格朗日子流形，并且相对 spin。

\par\smallskip\noindent\textup{(O.1)}\quad
对每个 $b_i\in\mathcal M^{\mathfrak b}(L_i;\Lambda_+)$，令
$\mathbf b_i=(\mathfrak b,b_i)$。则有弗洛尔上同调
\[
  HF\bigl((L_1,\mathbf b_1),(L_0,\mathbf b_0);\Lambda_0\bigr).
\]

\par\smallskip\noindent\textup{(O.2)}\quad
在 $(L_1,\mathbf b_1)=(L_0,\mathbf b_0)=(L,\mathbf b)$ 的情形，弗洛尔上同调群
\[
  HF\bigl((L,\mathbf b),(L,\mathbf b);\Lambda_0\bigr)
\]
具有环结构。

\par\smallskip\noindent\textup{(O.3)}\quad
弗洛尔上同调
$HF((L_1,\mathbf b_1),(L_0,\mathbf b_0);\Lambda_0)$
具有滤过双模结构；它的左作用环和右作用环分别为
\[
  HF\bigl((L_1,\mathbf b_1),(L_1,\mathbf b_1);\Lambda_0\bigr),
  \qquad
  HF\bigl((L_0,\mathbf b_0),(L_0,\mathbf b_0);\Lambda_0\bigr).
\]

\par\smallskip\noindent\textup{(O.4)}\quad
若对某个整数 $m$ 有 $m[\omega]\in H^2(M,L;\mathbb Z)$，则在上述陈述中，可以把系数环 $\Lambda_0$ 缩减为形式 Puiseux 级数环\liftIndex{I021}\liftIndex{I363}
$\mathbb Q[[T^{1/m}]]\subset\Lambda_0$。
\end{corollary}

\begin{remark}{1.17}\label{rem:1-17}
(1) 与定理 M 的情形不同，当马斯洛夫类为零时，为得到
$\mathcal M^{\mathfrak b}(L;\Lambda_+)$，不需要再把
$\mathfrak P^{\mathfrak b}_L$ 的零点集除以一个等价关系。参见第 5.4.6 小节的引理 5.4.71。

(2) 事实上，通过考虑拉格朗日子流形的任意有限链
$(L_0,L_1,\ldots,L_k)$，可以把本推论中的陈述扩张到滤过 $A_\infty$ 范畴\liftIndex{I006}的语境。本书只讨论到对象（一个拉格朗日子流形）和态射（两个拉格朗日子流形）的定义；关于 $A_\infty$ 范畴一般语境下的讨论，参见\textbf{[Fuk02II]}。
\end{remark}

推论 O 是定理 G 和定理 M 的一个小修改。只有注 1.17 (1) 中提到的要点以及陈述 (O.4)，不能立即由定理 G、M 得到。我们将在第 5.4.6 小节讨论这些要点并完成推论 O 的证明。

现在用形式概形的语言重述推论 O。考虑推论 O 的情形，令
\[
  m=\dim H^1(L;\mathbb Q),
  \qquad
  \ell=\dim H^2(L;\mathbb Q),
\]
并固定 $H^1(L;\mathbb Q)$ 的一组基。以 $s_1,\ldots,s_m$ 表示 $H^1(L;\mathbb Q)$ 上相应的坐标，并令
\[
  R(\mathfrak P^{\mathfrak b}_L;\Lambda_0)
  =\frac{\Lambda_0[[s_1,\ldots,s_m]]}
  {((\mathfrak P^{\mathfrak b}_L)_1,\ldots,
    (\mathfrak P^{\mathfrak b}_L)_\ell)}.
\]
这里把
$\mathfrak P^{\mathfrak b}_L=((\mathfrak P^{\mathfrak b}_L)_1,\ldots,
(\mathfrak P^{\mathfrak b}_L)_\ell)$
看作 $s_1,\ldots,s_m$ 的形式幂级数。

\begin{corollary}{P}\label{cor:intro-P}
设 $L,M$ 如推论 O。环 $R(\mathfrak P^{\mathfrak b}_L;\Lambda_0)$ 是 $L$ 的不变量。此外，对每个相对 spin 的对子 $(L_0,L_1)$，可以赋予一个上链复形
$\bigl(\mathfrak D(L_1,L_0;\Lambda_0),\mathfrak d\bigr)$；其各项是环对
\[
  \bigl(R(\mathfrak P^{\mathfrak b}_{L_1};\Lambda_0),
    R(\mathfrak P^{\mathfrak b}_{L_0};\Lambda_0)\bigr)
\]
上有限生成、连续且完备的双模，并满足如下性质：若
$b_i\in\mathcal M^{\mathfrak b}(L_i;\Lambda_+)$，且
$\varphi_{b_i}:R(\mathfrak P^{\mathfrak b}_{L_i};\Lambda_0)\to\Lambda_0$
是它的代表，则复形
\[
\begin{split}
\bigl(&{\Lambda_0}_{\varphi_{b_1}}
 \otimes_{R(\mathfrak P^{\mathfrak b}_{L_1};\Lambda_0)}
 \mathfrak D(L_1,L_0;\Lambda_0)
 \otimes_{R(\mathfrak P^{\mathfrak b}_{L_0};\Lambda_0)}
 {_{\varphi_{b_0}}\Lambda_0},\\
 &1\otimes\mathfrak d\otimes1\bigr)
\end{split}
\]
的上同调同构于弗洛尔上同调
\[
  HF\bigl((L_1,\mathbf b_1),(L_0,\mathbf b_0);\Lambda_0\bigr),
  \qquad \mathbf b_i=(\mathfrak b,b_i).
\]
在链同伦等价意义下，双模
$\bigl(\mathfrak D(L_1,L_0;\Lambda_0),\mathfrak d\bigr)$
与构造中所作的各种选择无关。陈述 (O.4) 仍然成立。
\end{corollary}

推论 P 是定理 N 的一个小修改，其证明见第 5.4.6 小节。

现在描述推论 O 和 P 中这些结构的猜想性镜像图景。除了要求 $M$ 为卡拉比--丘流形、$L$ 的马斯洛夫类为零之外，为了在代数而非超越的语境中描述这个镜像，我们还需要对 $M$ 和 $L$ 作一些限制。

\begin{statement}{条件}{1.18}\label{cond:1-18}
假设 $L$ 是有理的。也就是说，存在某个正整数\translatornote{原书写作“非负整数”，但 $m=0$ 会使该条件变成空条件，并使下文的 $T^{1/m}$ 无定义，故改为“正整数”。} $m$，使得对所有
$\beta\in H_2(M,L;\mathbb Z)$ 都有
\[
  (m\omega)\cap\beta\in\mathbb Z.
\]
\end{statement}

注意，这自动蕴含 $m[\omega]\in H^2(M,L;\mathbb Z)$。我们考虑 Puiseux 级数环 $\mathbb Q[[T^{1/m}]]\subset\Lambda_0$。（关于这一假设的更多内容，见\textbf{[Fuk03I]}。）

\begin{conjecture}{Q}\label{conj:intro-Q}
设 $(M,\omega)$ 是卡拉比--丘流形\liftIndex{I011}，并且对某个整数 $m$ 有
$m[\omega]\in H^2(M;\mathbb Z)$。考虑元素
\[
  \mathfrak b\in H^2(M;\mathbb Q[[T^{1/m}]])
  \subset H^2(M;\Lambda_0),
\]
并假设形式幂级数 $\mathfrak b$ 在 $T=0$ 处的收敛半径严格为正。则可赋予一个全纯族
\begin{equation}
  \pi:\mathfrak X_{(M,\omega,\mathfrak b)}\longrightarrow D^2(\epsilon)
  \tag{1.19}\label{eq:1-19}
\end{equation}
定义在某个充分小的 $\epsilon>0$ 所对应的圆盘 $D^2(\epsilon)\subset\mathbb C$ 上，使得对每个 $q\ne0$，纤维 $\pi^{-1}(q)$ 都是卡拉比--丘流形；$q=0$ 处的纤维 $\pi^{-1}(0)$ 是奇异的，并且当 $q\to0$ 时，这个族定义了一个极大退化族\liftIndex{I012}。
\end{conjecture}

极大退化族的定义见\textbf{[LTY05]}。下面简要解释，猜想 Q 如何由镜像对称研究者广泛相信但尚未证明的若干事实推出。（除了 $\mathfrak b$ 进入构造的方式之外，以下描述并不是新的；例如见\textbf{[Fuk03II]}。）

令
\[
  \mathfrak b=\sum_{k=1}^{\infty}\mathfrak b_kT^{k/m},
  \qquad \mathfrak b_k\in H^2(M;\mathbb Q).
\]
选取代表各个类 $\mathfrak b_k$ 的闭 $2$-形式，仍用同一符号表示它们。作出这些选择后，假设该级数的收敛半径为正。对
\[
  \tau\in\mathfrak h=\{\tau\in\mathbb C\mid\operatorname{Im}\tau>0\},
\]
令 $q=e^{2\pi\sqrt{-1}\tau/m}$。启发式地，可以把 $q$ 看作 $T^{1/m}$；不过 $T$ 是形式参数，而 $q\in D^2(\epsilon)$ 是真正的复数\translatornote{原书此处写作 $D^1(\epsilon)$；但式 (1.19) 已将底空间定义为复平面中的圆盘 $D^2(\epsilon)$，且 $q$ 为复数，故改正。}。对每个这样的 $\tau$，定义复值闭 $2$-形式
\begin{equation}
  \omega_\tau(\mathfrak b)
  =\frac{2\pi}{\sqrt{-1}}
   \left(\tau\omega+\sum_k\mathfrak b_kq^k\right).
  \tag{1.20}\label{eq:1-20}
\end{equation}
由 $\mathfrak b$ 的收敛性假设，\eqref{eq:1-20} 表明
$\omega_\tau(\mathfrak b)$ 随 $\tau\in\mathfrak h$ 全纯变化。也容易看出，只要
$\operatorname{Im}\tau$ 足够大，$\operatorname{Re}(\omega_\tau(\mathfrak b))$ 就非退化，因而在 $M$ 上定义一个辛形式。我们把族~\eqref{eq:1-20} 看作 $M$ 上复化辛结构的一个全纯族。$\omega_\tau(\mathfrak b)$ 的虚部称为 $B$-场\liftIndex{I007}。注意，$m[\omega]\in H^2(M;\mathbb Z)$ 蕴含
\begin{equation}
  \tau-\tau'\in m\mathbb Z
  \quad\Longrightarrow\quad
  \omega_\tau(\mathfrak b)-\omega_{\tau'}(\mathfrak b)
  \in2\pi\sqrt{-1}\,H^2(M;\mathbb Z).
  \tag{1.21}\label{eq:1-21}
\end{equation}

下面给出一个被广泛相信为真的事实：

\begin{quote}\itshape
对每个带 $B$-场的卡拉比--丘流形 $(M,\omega_\tau(\mathfrak b))$，都可以赋予同样是卡拉比--丘流形的镜像 $X_{\omega_\tau(\mathfrak b)}$。此外，如果两个复化辛结构之差属于
$2\pi\sqrt{-1}\,H^2(M;\mathbb Z)$，那么它们对应于同一个镜像流形。
\end{quote}

因此，由~\eqref{eq:1-21} 可以预期，我们已经定义了~\eqref{eq:1-19} 中提到的、以 $q$ 为参数的族。

再注意，对 $\beta\in H_2(M;\mathbb Z)$ 有恒等式
\begin{equation}
  \exp\bigl(-\omega_\tau(\mathfrak b)[\beta]\bigr)
  =q^{m(\omega[\beta])}
   \exp\left(2\pi\sqrt{-1}\sum_k(\mathfrak b_k[\beta])q^k\right).
  \tag{1.22}\label{eq:1-22}
\end{equation}
由~\eqref{eq:1-20} 和~\eqref{eq:1-22} 可见，当 $q\to0$ 时，上同调类 $\mathfrak b$ 以下述方式形变辛形式 $\omega$：
\[
  \frac{\sqrt{-1}\,\omega_\tau(\mathfrak b)}{2\pi\tau}
  \longrightarrow\omega
  \qquad\text{当 }\operatorname{Im}\tau\to\infty.
\]
这正是记号 $\mathcal M_{\mathrm{def}}(L)$ 中使用 ``def'' 的原因。

现在要表述一个猜想\liftIndex{I019}，用以描述对子 $(L,b)$ 的猜想性镜像对象；这里 $L$ 是满足条件 1.18 与 $\mu_L=0$ 的 spin 拉格朗日子流形，并且
$b\in\mathcal M^{\mathfrak b}(L;\Lambda_+)$。（稍微修改我们的构造，也可以把 $L$ 上的一个平坦复向量丛包括进来。这里略去讨论；关于这种修改的详细说明可参见\textbf{[Fuk02II]}。）

考虑 $\mathfrak X_{(M,\omega,\mathfrak b)}$ 沿 $\pi^{-1}(0)$ 的形式完备化
$\widehat{\mathfrak X}_{(M,\omega,\mathfrak b)}$（其含义见\textbf{[GrDi60]} I，第 10 节）。设 $i=1,2$ 时，$\hat{\mathfrak F}_i$ 是
$\widehat{\mathfrak X}_{(M,\omega,\mathfrak b)}$ 上凝聚层导出范畴\liftIndex{I013}中的对象。$\operatorname{Hom}$ 的右导出函子
$\operatorname{RHom}(\hat{\mathfrak F}_1,\hat{\mathfrak F}_2)$
定义为一个 $\mathbb C[[q]]$-模。（因为取了形式完备化，此时 $q$ 是形式参数。）

\begin{remark}{1.23}\label{rem:1-23}
(1) 奇异纤维 $\pi^{-1}(0)$ 似乎不存在典范选择。因此实际需要考虑
\[
  \operatorname{RHom}(\hat{\mathfrak F}_1,\hat{\mathfrak F}_2)
  \otimes_{\mathbb C[[q]]}\mathbb C[[q]][q^{-1}].
\]
更精确的表述见\textbf{[Fuk03I]} 第 4 节。从\textbf{[Ray74]} 的观点来看，在取 $\operatorname{RHom}$ 时忘掉 $0$ 处纤维的信息，意味着把 $\hat{\mathfrak F}_i$ 看作相应刚性解析空间的导出范畴中的对象。这一点可能与注 1.12 有关。

(2) 为使形式概形 $\widehat{\mathfrak X}_{(M,\omega,\mathfrak b)}$ 存在，我们甚至可能不需要在猜想 Q 中假设 $\mathfrak b$ 收敛。
\end{remark}

\begin{conjecture}{R}\label{conj:intro-R}
对每个对子 $(L,b)$，其中 $L$ 是满足条件 1.18 和 $\mu_L=0$ 的 spin 拉格朗日子流形，且
$b\in\mathcal M^{\mathfrak b}(L;\Lambda_+)$，都可以赋予
$\widehat{\mathfrak X}_{(M,\omega,\mathfrak b)}$ 上凝聚层导出范畴中的一个对象
$\hat{\mathfrak F}(L,b)$，并使下列结论成立。

\par\smallskip\noindent\textup{(R.1)}\quad
$\mathcal M^{\mathfrak b}(L;\Lambda_+)$ 中 $b$ 的形式邻域，同构于
$\hat{\mathfrak F}(L,b)$ 在其模空间中的一个形式邻域。

\par\smallskip\noindent\textup{(R.2)}\quad
存在典范同构
\[
\begin{split}
  &HF\bigl((L_1,\mathbf b_1),(L_0,\mathbf b_0);
    \mathbb C[[q]][q^{-1}]\bigr)\\
  &\qquad\cong
  \operatorname{RHom}\bigl(\hat{\mathfrak F}(L_1,b_1),
    \hat{\mathfrak F}(L_0,b_0)\bigr)
  \otimes_{\mathbb C[[q]]}\mathbb C[[q]][q^{-1}].
\end{split}
\]
这里 $\mathbf b_i=(\mathfrak b,b_i)$。（由 (O.4) 可知，左端的弗洛尔上同调确实可以在环 $\mathbb C[[q]][q^{-1}]$ 上定义。）

\par\smallskip\noindent\textup{(R.3)}\quad
(O.2)、(O.3) 中的环结构和双模结构，分别同构于 (R.2) 右端的相应结构。
\end{conjecture}

如\textbf{[Fuk02II]} 所述，通过考虑拉格朗日子流形的任意有限链，可以把猜想 R 推广到滤过 $A_\infty$ 范畴\liftIndex{I006}的语境。

\begin{remark}{1.24}\label{rem:1-24}
需要指出，导出范畴中对象的模空间，在代数几何里尚不是一个已经充分确立的概念。在许多重要情形中，镜像对象是
$\widehat{\mathfrak X}_{(M,\omega,\mathfrak b)}$ 上的凝聚层，而不只是其导出范畴中的一个对象。在这些情形下，(R.1) 中提到的模空间在代数几何里已有充分确立的定义。

对一般情形，可以利用本书研究的滤过 $A_\infty$ 代数概念，赋予 (R.1) 严格含义。先注意，(R.2) 中出现的两个上同调群，都是某个滤过 $A_\infty$ 双模的 $\delta_{\mathbf b_1,\mathbf b_0}$-上同调。对弗洛尔上同调，这由本书结果得到，更确切地说，由定理 A 和定义 3.6.9 得到；对 $\operatorname{RHom}$，则由定理 W 以及同调代数中 $\operatorname{RHom}$ 的标准构造得到。可以把 (R.3) 加强为一个猜想，即二者作为滤过 $A_\infty$ 双模同构。再回忆，模空间的形式邻域由（滤过）$A_\infty$ 结构控制；具体说，(1.8) 正是决定模空间形式邻域的方程（更多说明见\textbf{[Fuk03II]}）。从这个意义上说，(R.3) 的加强版本把 (R.1) 包含为一个特殊情形。
\end{remark}

注意，上面的 (R.1) 意味着，推论 O 所给的形式映射
$\mathfrak P^{\mathfrak b}_L$ 与其镜像上全纯向量丛（或凝聚层）模空间的仓西映射相同。作代换 $T^{1/m}=q$ 后，$\mathfrak P^{\mathfrak b}_L$ 可视为映射
\[
  H^1(L;q\mathbb C[[q]])\longrightarrow H^2(L;q\mathbb C[[q]])
\]
（见\textbf{[Fuk03I]}），并具有形式
\begin{equation}
  \mathfrak P^{\mathfrak b}_L(b)
  =\sum_{k=0}^{\infty}
    \mathfrak m_k(\underbrace{b,\ldots,b}_{k}).
  \tag{1.25}\label{eq:1-25}
\end{equation}
严格说来，\eqref{eq:1-25} 是链层上的公式。要在上同调
$H^1(L;\mathbb C[[q]])$ 上定义 $\mathfrak P^{\mathfrak b}_L$，需要使用后面的定理 W 以及第 5.4.3 小节定义的滤过 $A_\infty$ 同态。

\eqref{eq:1-25} 的右端可以看作以 $L$ 为边界的伪全纯圆盘数目的生成函数；它是格罗莫夫--威腾势的开弦类比。（格罗莫夫--威腾势是亏格零格罗莫夫--威腾不变量的生成函数；这些不变量通过计数伪全纯球面来定义。）（当 $\dim L=3$ 时，映射 $\mathfrak P^{\mathfrak b}_L$ 是某个超势的微分，因而是格罗莫夫--威腾势更直接的类比；见第 3.6.4 小节。）因此，(R.1) 是关于以 $L$ 为边界的伪全纯圆盘数目与其镜像上向量丛形变理论之关系的猜想。从这个意义上说，猜想 R 可以看作著名猜想
\[
  \text{量子杯积}\liftIndex{I025}=\text{汤川耦合}\liftIndex{I034}
\]
的开弦版本。这个猜想由菲利普·坎德拉斯（Philip Candelas）等人\textbf{[COGP91]} 发现，并由纪梵塔尔（Alexander Givental）\textbf{[Giv96]} 部分证明（另见\textbf{[LLY97]}）。这种开弦类比最先由孔采维奇在椭圆曲线情形发现\textbf{[Kon95I]}。Alexander Polishchuk--埃里克·扎斯洛（Eric Zaslow）\textbf{[PoZa98]} 进一步研究了椭圆曲线情形，第一署名作者又把它部分推广到高维\textbf{[Fuk02III]}。近来又有若干情形得到研究，其中包括第 1.3 节末提到的 (1.14.1) 和 (1.14.5)。

下面结合推论 P 与一些刚性解析几何，概述猜想 R 的一个证明方案。

\begin{remark}{1.26}\label{rem:1-26}
用刚性解析几何证明同调镜像对称的想法来自孔采维奇和 Yan Soibelman\textbf{[KoSo01]}。还有一种与此密切相关但有所不同的方法，即第一署名作者提出的渐近分析方法\textbf{[Fuk05I]}。另请注意，马克·格罗斯（Mark Gross）和贝恩德·西伯特（Bernd Siebert）在构造镜像流形的尝试中使用了对数概形\textbf{[GrSi03I,03II]}，其风格也颇为相近。
\end{remark}

为使叙述简单，本节余下部分只考虑 $\mathfrak b=0$ 的情形。

受斯特罗明格--丘成桐--扎斯洛（Strominger--Yau--Zaslow）图景\liftIndex{I030}\textbf{[SYZ96]} 启发，从一个 $2n$ 维辛流形 $M$ 出发。它在流形 $B$ 上有一个奇异纤维化，其一般纤维是马斯洛夫类为零的拉格朗日环面 $T^n$。假设 $\omega$ 是整的，即 $[\omega]\in H^2(M;\mathbb Z)$。以 $B_{\mathrm{sm}}$ 表示 $B$ 中所有使纤维 $F_y$ 光滑的点 $y$ 所组成的子集。在对该纤维化所作的假设下，可以在 $B_{\mathrm{sm}}$ 上构造一个
$GL(n;\mathbb Z)\widetilde{\times}\mathbb Z^n$（平坦仿射）结构（见\textbf{[Dui80, KoSo01, GrSi03I]}）。还假设
$\operatorname{codim}(B\setminus B_{\mathrm{sm}})\ge2$。定义 $B^{\mathbb Q}$ 为仿射坐标均为有理数的点所成的集合。当 $\pi_1(M)=1$ 时，$y\in B^{\mathbb Q}$ 等价于 $F_y$ 是有理拉格朗日子流形\liftIndex{I021}\liftIndex{I363}。

\begin{remark}{1.27}\label{rem:1-27}
设 $y\in B^{\mathbb Q}$。若 $F_y$ 是奇异的，本书没有给出 $F_y$ 的弗洛尔上同调之定义。在许多重要情形中，奇异纤维 $F_y$ 是浸入拉格朗日子流形。因此，很有必要把我们的理论扩张到这种情形（见\textbf{[Aka05, AkJo08]}）。
\end{remark}

\begin{conjecture}{S}\label{conj:intro-S}
设 $y\in B^{\mathbb Q}$，并令 $\mathfrak P_{F_y}$ 为~\eqref{eq:1-25} 中引入的形式幂级数映射。则有以下等式\translatornote{原书此处及下文两处均写作 $H^2(F_y;\Lambda_+)$；但 $\mathfrak P_{F_y}$ 的定义域是一次上同调中的边界链空间，且环面纤维的刚性点集为 $\Lambda_+^n$，故三处统一改为 $H^1$。}：
\[
  \mathfrak P_{F_y}\equiv0,
  \qquad\text{即}\qquad
  \mathcal M(F_y;\Lambda_+)=H^1(F_y;\Lambda_+).
\]
\end{conjecture}

现在考虑猜想 R 中的对子 $(L,b)$。由推论 P 得到
\begin{equation}
  \left(
    {\Lambda_0}_{\varphi_b}
    \mathbin{\widehat\otimes}_{R(\mathfrak P_L;\Lambda_0)}
    \mathfrak D(L,F_y;\Lambda_0),
    1\otimes\mathfrak d
  \right).
  \tag{1.28}\label{eq:1-28}
\end{equation}
它定义了形式概形
\[
  \operatorname{Spf}(R(\mathfrak P_{F_y};\Lambda_0))
  =\operatorname{Spf}(\Lambda_0[[s_1,\ldots,s_n]])
\]
上凝聚层导出范畴\liftIndex{I013}中的一个对象。回忆，形式概形
$\operatorname{Spf}(R(\mathfrak P_{F_y};\Lambda_0))$ 的刚性点集是
$\Lambda_+^n$，其中 $s_1,\ldots,s_n$ 是它的坐标。事实上，自然坐标不是 $s_i$，而是它们的指数
$e^{s_i}=x_i$（对此可见第 3.6.4 小节）。集合
$\exp(H^1(F_y;\Lambda_+))$ 是
$(1,\ldots,1)\in(\Lambda_0\setminus\{0\})^n$ 的一个无穷小（开）邻域。

\begin{conjecture}{T}\label{conj:intro-T}
\eqref{eq:1-28} 中给出的对象可下降为开集
$\exp(H^1(F_y;\Lambda_+))$ 上的对象，并可扩张成环形区域
\[
  A_{1-\epsilon,1+\epsilon}(x_1,\ldots,x_n)
  =\left\{(x_1,\ldots,x_n)\in(\Lambda_0\setminus\{0\})^n
  \mathrel{\Big|}1-\epsilon<\lVert x_i\rVert<1+\epsilon\right\}
\]
上的对象，其中 $\epsilon>0$ 充分小。这里 $\lVert\mathord\cdot\rVert$ 是如下定义的非阿基米德范数：
\[
  \left\lVert\sum_i a_iT^{\lambda_i}\right\rVert=e^{-\lambda_1},
\]
其中 $\lambda_1$ 是首指数，即 $a_1\ne0$ 且当 $i>1$ 时 $\lambda_i>\lambda_1$。
\end{conjecture}

为了在上述无穷小邻域及其坐标的讨论中指明基点 $y\in B$，用 $x_i^y$ 代替 $x_i$。也就是说，定义 $x_i^y=\exp(s_i)$，其中 $s_i$ 是为
$\operatorname{Spf}(R(\mathfrak P_{F_y};\Lambda_0))$ 选定的坐标。设
$y'=(y'_1,\ldots,y'_n)$ 是 $y=(y_1,\ldots,y_n)$ 的邻域中的一点，二者均用 $B$ 上给定的仿射坐标表示。比较分别以拉格朗日子流形 $F_y$ 与 $F_{y'}$ 为边界的圆盘之辛面积，可以看到，按下式定义坐标 $x_i^{y'}$ 与 $x_i^y$ 的关系是自然的：
\begin{equation}
  x_i^{y'}=T^{y_i'-y_i}x_i^y.
  \tag{1.29}\label{eq:1-29}
\end{equation}

\begin{conjecture}{U}\label{conj:intro-U}
存在（刚性解析空间的）坐标变换映射 $\Phi_{yy'}$，从
\begin{equation}
\begin{split}
  &A_{1-\epsilon,1+\epsilon}(x_1^y,\ldots,x_n^y)
  \cap A_{1-\epsilon,1+\epsilon}(x_1^{y'},\ldots,x_n^{y'})\\
  &\hspace{9em}\subset
  A_{1-\epsilon,1+\epsilon}(x_1^y,\ldots,x_n^y)
\end{split}
  \tag{1.30}\label{eq:1-30}
\end{equation}
映到
\begin{equation}
\begin{split}
  &A_{1-\epsilon,1+\epsilon}(x_1^y,\ldots,x_n^y)
  \cap A_{1-\epsilon,1+\epsilon}(x_1^{y'},\ldots,x_n^{y'})\\
  &\hspace{9em}\subset
  A_{1-\epsilon,1+\epsilon}(x_1^{y'},\ldots,x_n^{y'}).
\end{split}
  \tag{1.31}\label{eq:1-31}
\end{equation}
映射 $\Phi_{yy'}$ 由一个适当的滤过 $A_\infty$ 同态
\[
  H(F_{y'};\Lambda_0)\longrightarrow H(F_y;\Lambda_0)
\]
诱导。可以用从~\eqref{eq:1-30} 到~\eqref{eq:1-31} 的转移映射
$\Phi_{yy'}$ 粘合
$\{A_{1-\epsilon,1+\epsilon}(x_1^y,\ldots,x_n^y)\}$，得到一个刚性解析空间 $\widehat{\mathfrak X}$。

此外，可以用 $\Phi^L_{yy'}$ 粘合猜想 T 中得到的对象，从而得到
$\widehat{\mathfrak X}$ 上凝聚层导出范畴中的对象
$\hat{\mathfrak F}(L,b)$。转移映射 $\Phi^L_{yy'}$ 由适当的滤过
$A_\infty$ 双模同态诱导。
\end{conjecture}

如果得到证明，猜想 U 就给出了构造猜想 R 中镜像对象
$\hat{\mathfrak F}(L,b)$ 的一种方法。于是该构造的函子性将蕴含：存在从 (R.2) 左端到其右端的滤过 $A_\infty$ 同态。要证明它是同构，还需要更多论证。

要实施上述方案，仍有其他困难。一个突出的难题，是研究奇异纤维及与之相关的各种结构。我们预期\textbf{[FOOO09II]} 的结果会对解决这一问题有用。

还要指出，上述方案使用了弗洛尔上同调的族版本。用弗洛尔上同调的族版本证明同调镜像对称的想法，可以追溯到 1997 年夏第一署名作者与孔采维奇在 IHES 的讨论。对阿贝尔簇的情形，这个想法已在\textbf{[Fuk02III]} 中得到部分实现。在这个特殊情形，\textbf{[Fuk02III]} 中出现的所有运算 $\mathfrak m_k$ 确实都收敛，因为它们可归结为多重 theta 函数；所以在\textbf{[Fuk02III]} 中，不需要把论述转换成形式概形或刚性解析空间的语言。另见\textbf{[Fuk02I]}。

现阶段，上述以弗洛尔上同调的族版本研究同调镜像对称猜想的方法，看来是最有希望的方法；我们希望能在不久的将来完成它。

结束本节之前，我们再指出这个理论的几种可能推广。本书对弗洛尔上同调和 $A_\infty$ 结构的构造，使用了从圆盘 $D^2$ 出发的伪全纯映射之模空间。也可以考虑从高亏格带边黎曼曲面出发的伪全纯映射，并在 $L$ 上考虑比平坦线丛更一般的平坦丛。如果对一个 $3$-流形余切丛的零截面施行这些构造，相应的开弦理论预期与陈--西蒙斯（Chern--Simons）扰动理论一致\textbf{[AxSi91I,91II, Bar91, GMM89, Kon94]}。威腾在\textbf{[Wit95]} 中指出了这一事实（另见\textbf{[Fuk96I,97III, FuOh97]}）。它的镜像将是 Bershadsky--Cecotti--Ooguri--Vafa 量子 Kodaira--Spencer 理论\textbf{[BCOV94]} 的一个版本；更确切地说，是该理论在凝聚层模空间上的类比。\textbf{[BCOV94]} 描述了流形复结构模空间上的量子 Kodaira--Spencer 理论。遗憾的是，要像猜想 R 那样，以数学上精确的形式表述这种一致性，仍超出我们目前所能达到的范围。
\enlargethispage{2\baselineskip}
不过，对亏格一的情形，也许可以得到某种严格结论。在辛几何一侧，加入量子贡献后，亏格一部分应对应于 Michael Hutchings 和 Yi-Jen Lee 定义的一个不变量\textbf{[HuLe99]}。这个不变量与伪全纯圆环的数目及 Reidemeister 扭率有关（另见\textbf{[Fuk97I, Lee01]}）。在复几何一侧，亏格一部分应对应于 $\bar\partial$ 算子的解析扭率，其系数取为
$\operatorname{Hom}(\hat{\mathfrak F}(L,b),\hat{\mathfrak F}(L,b))$。$\bar\partial$ 算子的解析扭率与凝聚层模空间上的 Quillen 度量等对象之间的关系，已有许多人讨论，例如尚-米歇尔·比斯姆（Jean-Michel Bismut）\textbf{[Bis98]}。研究猜想 R 的高亏格类比时，第一步应当是证明亏格一情形下定义的这些不变量彼此相同。
